\documentclass[../DoAn.tex]{subfiles}
\begin{document}
\sloppy

% ===== ĐỊNH NGHĨA TỪ KHÓA THUẬT TOÁN CHO CHƯƠNG 5 =====
\SetKwIF{If}{ElseIf}{Else}{IF}{THEN}{ELSE IF}{ELSE}{END IF}
\SetKwFor{For}{FOR}{DO}{END FOR}
\SetKwFor{While}{WHILE}{DO}{END WHILE}
\SetKwSwitch{Switch}{Case}{Other}{CASE}{OF}{}{}{END CASE}{}
\SetKw{Begin}{BEGIN}
\SetKw{End}{END}
\SetKw{KwReturn}{RETURN}
\SetKw{KwBreak}{BREAK}
\SetKw{KwCall}{CALL}
\SetKw{KwException}{EXCEPTION}
\SetAlgoNoLine  % Bỏ đường kẻ ngang trong tất cả thuật toán
\SetAlgoSkip{}          % bỏ skip trên/dưới algorithm
\setlength{\algomargin}{1em}  % giảm indent margin
% =======================================================

Bốn chương trước là phần tôi cố gắng viết chặt chẽ theo quy chuẩn học thuật
--- bài toán, yêu cầu, công nghệ, thiết kế. Chương này là phần tôi muốn kể
thật hơn: những lúc tưởng chừng bế tắc, những quyết định thiết kế phải làm
đi làm lại, và lý do tại sao cuối cùng chọn hướng đó. Ba đóng góp được trình
bày không phải theo kiểu ``đây là giải pháp đúng'' mà theo kiểu ``đây là hành
trình tìm ra giải pháp đó''.

Mỗi đóng góp vẫn trả lời bốn câu hỏi cốt lõi: \textit{(i)} tính năng giải
quyết vấn đề gì; \textit{(ii)} đối tượng nào được hưởng lợi; \textit{(iii)}
điểm khác biệt so với giải pháp hiện có là gì; \textit{(iv)} giá trị mang lại
cho hệ thống và người dùng ra sao. Chương được tổ chức thành bốn mục: giải
pháp trực quan hóa không gian phòng ký túc xá~(5.1), giải pháp phân bổ phòng
tự động dựa trên quota~(5.2), giải pháp kiến trúc phân bổ linh hoạt theo năm
học~(5.3) và tổng kết chương~(5.4).

% =====================================================================
\section{Giải pháp trực quan hóa không gian phòng ký túc xá}
\label{sec:contrib_visualization}
% =====================================================================

\subsection{Dẫn dắt và bài toán}
\label{subsec:viz_problem}

\subsubsection{Vấn đề cần giải quyết}

Phần lớn sinh viên nhập học tại HUST đến từ các tỉnh thành ngoài Hà Nội và
chưa từng đặt chân đến ký túc xá trước khi đăng ký. Hệ thống hiện tại chỉ
cung cấp ảnh tĩnh và mô tả văn bản --- không đủ để hình thành nhận thức không
gian thực tế, vì mắt người cần tự do chuyển hướng nhìn mới cảm nhận được
chiều sâu và tỷ lệ. Hệ quả là một tỷ lệ đáng kể sinh viên phát hiện phòng
không khớp kỳ vọng sau khi nhận, dẫn đến yêu cầu đổi phòng và tăng tải công
việc cho cán bộ quản lý.

Bài toán đặt ra là: \textit{thiết kế module trực quan hóa không gian đủ chân
thực để sinh viên ra quyết định đăng ký từ xa, đảm bảo hiệu suất tải hợp lý
trên web và khả dụng trên cả thiết bị di động.}

\subsubsection{Đối tượng hưởng lợi}

\textit{Sinh viên đăng ký} có thể ``đi vào'' phòng, nhìn quanh 360°, di
chuyển tự do trong không gian 3D trước khi nhấn nút đăng ký --- quyết định
dựa trên trải nghiệm thực tế thay vì suy đoán. \textit{Cán bộ quản lý} giảm
số yêu cầu đổi phòng sau đăng ký và có công cụ tự cập nhật nội dung trực quan
mà không cần lập trình viên. \textit{Hệ thống eDorm} hưởng lợi qua chất lượng
dữ liệu: lựa chọn phòng của sinh viên phản ánh đúng sở thích hơn, giảm nhiễu
cho thuật toán phân bổ.

\subsubsection{Điểm khác biệt so với giải pháp hiện có}

Hầu hết ký túc xá đại học tại Việt Nam dừng lại ở ảnh tĩnh hoặc video tour
tuyến tính --- cả hai đều chỉ cho phép xem thụ động. Điểm khác biệt then chốt
của giải pháp đề xuất là \textit{chuyển dịch từ xem thụ động sang trải nghiệm
chủ động}: người dùng tự quyết định nhìn vào đâu và đứng ở đâu trong phòng.
Bảng~\ref{tab:viz_comparison} tổng hợp so sánh ba phương án.

\begin{table}[H]
\centering
\caption{So sánh các phương án trực quan hóa không gian phòng}
\label{tab:viz_comparison}
\small
\begin{tabular}{|p{2.8cm}|p{2.6cm}|p{2.6cm}|p{5.2cm}|}
\hline
\textbf{Tiêu chí} & \textbf{Ảnh tĩnh} & \textbf{Video tour} & \textbf{Giải pháp đề xuất (360° + 3D)} \\
\hline
Nhận thức không gian & Rất hạn chế & Trung bình & Cao --- cảm nhận chiều sâu và tỷ lệ thực \\
\hline
Tính tương tác & Không có & Thụ động & Chủ động --- người dùng tự chọn góc nhìn \\
\hline
Di chuyển tự do & Không & Không & Có (chế độ Free Explore trên desktop) \\
\hline
Tương thích thiết bị & Mọi thiết bị & Mọi thiết bị & 360° mọi thiết bị; 3D trên desktop \\
\hline
Cập nhật nội dung & Cần kỹ thuật & Cần kỹ thuật & Cán bộ tự cập nhật qua giao diện \\
\hline
Thời gian sẵn sàng & Ngay lập tức & Phụ thuộc băng thông & < 0,8 giây đến tương tác đầu tiên \\
\hline
\end{tabular}
\end{table}

\subsection{Giải pháp}
\label{subsec:viz_solution}

\subsubsection{Hành trình thử nghiệm --- từ ảnh tĩnh đến mô hình 3D}

Đây là đóng góp mà tôi phải thừa nhận đã tốn nhiều công sức và vòng lặp
thử-sai nhất trong toàn bộ quá trình làm đồ án. Không có một giải pháp nào
hiện ra ngay từ đầu --- mỗi lần tưởng chừng đã tìm được hướng đi đúng, lại
gặp một hạn chế mới buộc phải điều chỉnh. Hành trình đó trải qua bốn giai
đoạn, mỗi giai đoạn giải quyết một hạn chế của giai đoạn trước, được minh
họa trong Hình~\ref{fig:viz_journey}.

\begin{figure}[H]
\centering
\includegraphics[width=0.98\textwidth]{Hinhve/quy_trinh_xay_dung_3d.png}
\caption{Hành trình bốn giai đoạn xây dựng giải pháp trực quan hóa}
\label{fig:viz_journey}
\end{figure}

\textbf{Giai đoạn 1 --- Ảnh tĩnh nhiều góc độ.}
Hướng tiếp cận đầu tiên là tự nhiên nhất: chụp thêm nhiều ảnh từ các góc độ
khác nhau trong phòng, kèm chú thích kích thước từng khu vực. Ý tưởng là nếu
có đủ ảnh bao phủ mọi góc, người dùng sẽ có đủ thông tin để hình dung. Tuy
nhiên, sau khi triển khai và thử nghiệm thực tế, kết quả không như kỳ vọng.
Vấn đề không phải là thiếu ảnh mà là bản chất của ảnh tĩnh: mắt người nhận
thức chiều sâu thông qua chuyển động và góc nhìn liên tục thay đổi --- điều
mà tập hợp ảnh rời rạc không thể cung cấp. Giai đoạn này cho thấy vấn đề
không phải là số lượng thông tin mà là \textit{loại} thông tin cần cung cấp.

\textbf{Giai đoạn 2 --- Ảnh equirectangular 360° và hệ thống hotspot.}
Nhận ra rằng cần cho người dùng khả năng tự chuyển hướng nhìn, hướng tiếp cận
thứ hai là chụp ảnh equirectangular 360° tại các vị trí đại diện trong phòng
(cạnh cửa vào, giữa phòng, góc gần cửa sổ) và triển khai thư viện Photo Sphere
Viewer~\cite{psv2023}. Đây là bước đột phá thực sự đầu tiên: lần đầu tiên
người dùng có thể kéo chuột và nhìn quanh toàn bộ không gian từ một điểm đứng
cố định. Cảm giác chiều sâu và tỷ lệ không gian bắt đầu hình thành rõ ràng
hơn nhiều so với ảnh tĩnh.

Tuy nhiên, hạn chế nhanh chóng bộc lộ: người dùng bị ``giam cầm'' tại một
điểm quan sát, không thể di chuyển sang vị trí khác. Một phòng ký túc xá có
nhiều khu vực cần quan sát (khu giường, khu học, khu để đồ, khu vệ sinh) ---
mỗi khu cần một vị trí quan sát riêng. Giải pháp là xây dựng hệ thống
\textit{hotspot điều hướng}: các điểm tương tác nổi lên trong ảnh 360°, cho
phép người dùng nhấn vào để dịch chuyển sang góc quan sát tiếp theo.

Việc đặt hotspot đúng vị trí trong không gian ảnh cầu là thách thức kỹ thuật
không nhỏ. Mỗi hotspot phải được biểu diễn bằng tọa độ góc $(\theta, \phi)$
trong hệ tọa độ cầu, trong khi vị trí vật lý của điểm đích lại được lưu trong
hệ tọa độ Descartes $(x, y, z)$ của Three.js. Công thức chuyển đổi:
\begin{equation}
\theta = \arctan\!\left(\frac{x}{-z}\right), \quad
\phi   = \arctan\!\left(\frac{y}{\sqrt{x^2+z^2}}\right)
\label{eq:cart_to_sphere}
\end{equation}
Khi triển khai thực tế, phát sinh thêm một vấn đề thứ cấp: làm sao để cán bộ
kỹ thuật (không phải lập trình viên) có thể tự định vị hotspot cho từng phòng
mới mà không cần can thiệp vào mã nguồn? Giải pháp là xây dựng một công cụ ghi
tọa độ tích hợp trực tiếp trong giao diện quản trị: cán bộ di chuyển con trỏ
đến vị trí mong muốn trong ảnh 360°, nhấn ``Lấy tọa độ'', và hệ thống tự động
áp dụng công thức~\eqref{eq:cart_to_sphere} rồi lưu $(\theta, \phi)$ vào cơ
sở dữ liệu. Đây là quyết định thiết kế quan trọng: giải quyết không chỉ bài
toán kỹ thuật mà còn bài toán vận hành dài hạn.

\textbf{Giai đoạn 3 --- Mô hình 3D tương tác và cuộc khủng hoảng kích thước file.}
Chế độ 360° với hotspot đã cải thiện đáng kể trải nghiệm, nhưng vẫn còn một
hạn chế nền tảng: người dùng chỉ có thể quan sát từ các vị trí đã được cán bộ
định nghĩa trước. Mục tiêu kế tiếp là cho phép di chuyển hoàn toàn tự do ---
người dùng tự chọn đứng ở đâu, nhìn vào đâu, giống như thực sự bước vào phòng.

Không có sẵn mô hình 3D nào cho các phòng ký túc xá HUST, không có file CAD,
không có bản vẽ kỹ thuật dạng số --- chỉ có bản vẽ giấy và thước dây. Tôi
phải tự học Blender~\cite{blender2024} từ đầu trong khi vẫn song song thực
hiện các phần khác của đồ án. Quá trình dựng mô hình bao gồm: đo đạc thực tế
từng kích thước trong phòng, dựng hình học cơ bản (tường, sàn, trần, cửa sổ,
cửa ra vào), mô hình hóa từng món đồ nội thất (giường tầng, bàn học, tủ cá
nhân, kệ sách), và thiết lập vật liệu cho từng bề mặt.

Sau vài tuần làm việc, mô hình đầu tiên hoàn thành. Và rồi vấn đề thực sự
xuất hiện: file GLB xuất ra từ Blender có kích thước hơn \textbf{200~MB}. Đây
là con số không tưởng cho một ứng dụng web --- với kết nối thông thường, người
dùng phải chờ hàng phút chỉ để tải về mô hình phòng, chưa kể bộ nhớ GPU cần
để render. Toàn bộ hướng tiếp cận tưởng chừng như phải từ bỏ. Thú thật là lúc
nhìn thấy con số đó, tôi đã nghĩ đến việc bỏ hẳn hướng 3D và quay về chỉ dùng
360°. Nếu không có ảnh 360° đã hoạt động tốt ở giai đoạn trước làm fallback,
có lẽ tôi đã dừng lại ở đây.

\textbf{Giai đoạn 4 --- Hành trình tối ưu hóa từ 200~MB về dưới 8~MB.}
Thay vì từ bỏ, tôi quyết định tìm hiểu tại sao file lại nặng đến vậy và có
thể làm gì để giảm. Điều quan trọng ở đây không phải là thử từng kỹ thuật một
cách ngẫu nhiên, mà trước tiên phải hiểu đúng nguồn gốc của con số 200~MB đến
từ đâu --- chỉ khi biết rõ \textit{tại sao} nặng mới có thể chọn đúng công cụ
để giảm. Quá trình phân tích cho thấy ba nguồn gốc chính: \textit{(i)} số lượng
đa giác quá lớn (mô hình ban đầu có hơn 400.000 đa giác vì các bề mặt phẳng
như sàn và tường được chia thành quá nhiều tam giác nhỏ); \textit{(ii)} vật
liệu PBR (Physically Based Rendering) với nhiều texture map độ phân giải cao
cho mỗi bề mặt; \textit{(iii)} dữ liệu hình học không được nén.

\textit{Decimate geometry:} Blender cung cấp modifier Decimate cho phép tự
động hợp nhất các đa giác trên bề mặt phẳng. Áp dụng Decimate aggressively
cho các bề mặt lớn (sàn, tường, trần) nhưng giữ nguyên cho đồ nội thất có
chi tiết, số đa giác giảm từ hơn 400.000 xuống khoảng 120.000 --- bước tiến
đáng kể nhưng chưa đủ.

\textit{Texture baking:} Thay vì giữ nhiều texture map riêng lẻ (albedo,
roughness, metallic, normal map) cho mỗi vật liệu, toàn bộ thông tin vật liệu
được ``nướng'' (bake) vào một texture map duy nhất kích thước 1024×1024~px cho
mỗi nhóm đối tượng. Kỹ thuật này giảm mạnh số lượng draw call trong quá trình
render và tổng kích thước dữ liệu texture. Sau bước này, kích thước file giảm
xuống còn khoảng 30--40~MB.

\textit{Nén Draco:} Thư viện Draco~\cite{google2023draco} của Google cung cấp
thuật toán nén hình học 3D chuyên dụng, giảm 40--60\% kích thước dữ liệu
vertex và face mà không làm thay đổi hình dạng mô hình. Sau khi áp dụng Draco
với mức nén cân bằng (để giữ chất lượng hình ảnh), file GLB cuối cùng đạt dưới
8~MB --- giảm hơn 25 lần so với mô hình ban đầu.

Hành trình tối ưu hóa được tổng kết trong Bảng~\ref{tab:glb_optimization}.

\begin{table}[H]
\centering
\caption{Tiến trình tối ưu hóa kích thước file GLB}
\label{tab:glb_optimization}
\small
\begin{tabular}{|p{4.2cm}|p{2.2cm}|p{2.5cm}|p{4.3cm}|}
\hline
\textbf{Giai đoạn} & \textbf{Kích thước} & \textbf{Số đa giác} & \textbf{Kỹ thuật áp dụng} \\
\hline
Mô hình gốc (chưa tối ưu) & $\sim$200~MB & $>$400.000 & --- \\
\hline
Sau Decimate & $\sim$35~MB & 120.000 & Hợp nhất đa giác bề mặt phẳng \\
\hline
Sau Texture Baking & $\sim$30~MB & 120.000 & Gộp vật liệu thành texture map \\
\hline
Sau nén Draco & < 8~MB & 22.000--25.000 & Nén hình học + geometry index \\
\hline
\end{tabular}
\end{table}

Ngoài vấn đề kích thước file, còn một thách thức chỉ lộ ra khi cầm điện thoại
lên và thực sự mở trang: trình duyệt không crash ngay --- nó load được nhưng
sau vài giây render thì đóng băng hoàn toàn. Tra log Chrome DevTools mobile
mới phát hiện lỗi thiếu WebGL extension \texttt{OES\allowbreak\_element\allowbreak\_index\allowbreak\_uint}
mà Three.js cần cho index buffer 32-bit khi render mô hình lớn. Thay vì cố
vượt qua giới hạn phần cứng, quyết định thiết kế được đưa ra: ẩn hoàn toàn
chế độ Free Explore trên thiết bị di động và chỉ cung cấp Quick Tour (360°).
Đây là đánh đổi có chủ ý --- tốt hơn là cung cấp trải nghiệm hạn chế nhưng
ổn định, hơn là cố render mô hình và gây crash trình duyệt.

\subsubsection{Kiến trúc hai chế độ song song}

Toàn bộ hành trình bốn giai đoạn hội tụ vào kiến trúc cuối cùng: hai chế độ
song song chia sẻ cùng một tầng dữ liệu, được trình bày trong
Hình~\ref{fig:visualization_architecture}.

\begin{figure}[H]
\centering
\includegraphics[width=0.8\textwidth]{Hinhve/kien_truc_truc_quan_hoa_phong.png}
\caption{Kiến trúc giải pháp trực quan hóa không gian phòng ký túc xá}
\label{fig:visualization_architecture}
\end{figure}

Module tổ chức thành hai nhánh chạy song song. Nhánh \textit{Quick Tour} xử
lý ảnh equirectangular 360° qua Photo Sphere Viewer, khởi tạo ngay lập tức
khi trang tải và khả dụng trên mọi thiết bị. Nhánh \textit{Free Explore} tải
mô hình GLB qua Three.js và chỉ được kích hoạt trên desktop khi mạng đủ nhanh.
Cả hai nhánh chia sẻ tầng siêu dữ liệu phòng và hệ thống hotspot, đảm bảo
tính nhất quán dữ liệu giữa hai chế độ.

Chiến lược tải song song là điểm quan trọng về hiệu suất: thay vì chờ GLB
tải xong mới hiển thị gì đó, Quick Tour được khởi tạo ngay lập tức trong khi
GLB tải ngầm --- giảm thời gian đến tương tác đầu tiên từ hàng giây xuống
dưới 0,8 giây. Bảng~\ref{tab:viz_modes} tóm tắt đặc điểm từng chế độ.

\begin{table}[H]
\centering
\caption{So sánh hai chế độ trực quan hóa}
\label{tab:viz_modes}
\small
\begin{tabular}{|p{2.8cm}|p{5.4cm}|p{5.0cm}|}
\hline
\textbf{Tiêu chí} & \textbf{Quick Tour (360°)} & \textbf{Free Explore (3D)} \\
\hline
Công nghệ & Photo Sphere Viewer + Three.js & Three.js WebGL + GLB \\
\hline
Điều hướng & Hotspot tọa độ cầu xác định trước & Di chuyển tự do WASD + chuột \\
\hline
Tải trọng & Thấp (JPEG tối ưu) & < 8~MB (GLB nén Draco) \\
\hline
Thiết bị hỗ trợ & Desktop + mobile + tablet & Desktop (WebGL đầy đủ) \\
\hline
Thời gian sẵn sàng & Ngay lập tức & Sau 1,2--1,8 giây (tải ngầm) \\
\hline
\end{tabular}
\end{table}

\subsubsection{Thuật toán tối ưu hóa mô hình 3D}

\begin{minipage}{\linewidth}
\paragraph{Mục tiêu.}
Giảm kích thước file GLB từ mô hình gốc (hàng trăm MB) xuống mức chấp nhận
được cho ứng dụng web (dưới 10~MB) mà không làm mất chất lượng hình ảnh cảm
nhận được bởi người dùng, thông qua ba kỹ thuật áp dụng tuần tự.

\begin{algorithm}[H]
\small
\caption{\textit{GLBModelOptimizationPipeline}}
\KwData{$RawModel$: mô hình Blender gốc; $TargetSize$: kích thước mục tiêu (MB);
        $QualityThreshold$: ngưỡng chất lượng chấp nhận được}
\KwResult{$OptimizedGLB$: file GLB đã tối ưu; $Metrics$: báo cáo kích thước và chất lượng}
\tcp{Bước 1: Decimate}
\For{$obj \in RawModel.objects$}{
  \If{$obj.type \in \{\text{FLOOR, WALL, CEILING}\}$}{
    $obj.decimateRatio \leftarrow 0.05$\; \tcp*[r]{Bề mặt phẳng --- giảm mạnh}
  }
  \Else{
    $obj.decimateRatio \leftarrow 0.4$\; \tcp*[r]{Đồ nội thất --- giữ chi tiết}
  }
  $\text{ApplyDecimate}(obj,\ obj.decimateRatio)$\;
}
\tcp{Bước 2: Texture Baking}
\For{$group \in \text{GroupByMaterial}(RawModel.objects)$}{
  $texMap \leftarrow \text{BakeToTexture}(group,\ \text{resolution:}\ 1024\times1024)$\;
  $\text{ReplaceMaterials}(group,\ texMap)$\;
}
\tcp{Bước 3: Nén Draco --- tìm mức nén tối ưu}
$rawGLB \leftarrow \text{ExportGLB}(RawModel)$\;
$level \leftarrow 7$\;
\While{$rawGLB.size > TargetSize$ \textbf{AND} $level \leq 10$}{
  $candidate \leftarrow \text{ApplyDraco}(rawGLB,\ \text{compressionLevel:}\ level)$\;
  \If{$\text{VisualQuality}(candidate) \geq QualityThreshold$}{
    $OptimizedGLB \leftarrow candidate$\;
    $level \leftarrow level + 1$\;
  }
  \Else{
    \KwBreak\; \tcp*[r]{Vượt ngưỡng chất lượng --- dừng tăng nén}
  }
}
$Metrics \leftarrow \{size{:}\ OptimizedGLB.size,\ polygons{:}\ \text{CountPolygons}(OptimizedGLB)\}$\;
\KwReturn $OptimizedGLB$, $Metrics$\;
\end{algorithm}
\end{minipage}

\subsubsection{Thuật toán khởi tạo và điều phối module}

\begin{minipage}{\linewidth}
\paragraph{Mục tiêu.}
Điều phối quá trình tải tài nguyên song song và phân nhánh theo phần cứng
thiết bị, đảm bảo người dùng luôn có tương tác đầu tiên trong dưới 0,8 giây
và không xảy ra crash trên GPU di động.

\begin{algorithm}[H]
\small
\caption{\textit{InteractiveDormitoryVisualization}}
\KwData{$RoomID$: mã phòng; $DeviceType$: loại thiết bị; $NetworkSpeed$: tốc độ mạng}
\KwResult{$Scene$: scene tương tác sẵn sàng; $Status$: trạng thái tải}
$metadata \leftarrow \text{LoadRoomMetadata}(RoomID)$\;
\If{$metadata = \varnothing$}{
  \KwException ``Phòng không tồn tại hoặc không có dữ liệu''\;
}
\tcp{Ưu tiên 1: Khởi tạo Quick Tour ngay lập tức cho mọi thiết bị}
$\text{InitQuickTour}(metadata.sphericalImages,\ metadata.hotspots)$\;
\If{$DeviceType \neq \text{MOBILE}$}{
  \tcp{Ưu tiên 2: Tải GLB ngầm song song, không chặn Quick Tour}
  \If{$NetworkSpeed \geq \text{THRESHOLD\_FAST}$}{
    $\text{LoadGLBAsync}(metadata.glbPath,\ \text{onReady: EnableFreeExplore})$\;
  }
}
\While{người dùng đang tương tác}{
  \If{$action = \text{ROTATE}$}{
    $\text{UpdateCamera}(delta)$\;
  }
  \ElseIf{$action = \text{HOTSPOT\_CLICK}$}{
    $\text{NavigateTo}(hotspotId)$\;
  }
  \ElseIf{$action = \text{SWITCH\_MODE}$}{
    $\text{ToggleMode}()$\;
  }
}
\KwReturn $Scene$, $Status \leftarrow \text{READY}$\;
\end{algorithm}
\end{minipage}

\medskip
\paragraph{Phân tích độ phức tạp.}
Giai đoạn khởi tạo có độ phức tạp $O(h)$ với $h$ là số hotspot cần đăng ký;
vòng lặp tương tác xử lý mỗi sự kiện trong $O(1)$. Chiến lược tải song song
giảm thời gian đến tương tác đầu tiên 3--5 lần so với tải tuần tự.

\subsection{Kết quả đạt được}
\label{subsec:viz_result}

\begin{table}[H]
\centering
\caption{Chỉ số kỹ thuật sau toàn bộ quá trình tối ưu hóa}
\label{tab:viz_metrics}
\small
\begin{tabular}{|p{5.8cm}|p{2.6cm}|p{4.4cm}|}
\hline
\textbf{Chỉ số} & \textbf{Giá trị} & \textbf{Ghi chú} \\
\hline
Kích thước GLB ban đầu & $\sim$200~MB & Mô hình Blender gốc, chưa tối ưu \\
\hline
Kích thước GLB sau tối ưu & < 8~MB & Sau Decimate + Baking + Draco \\
\hline
Hệ số nén tổng thể & $>$25$\times$ & Không mất chất lượng cảm nhận \\
\hline
Số đa giác sau tối ưu & 22.000--25.000 & Giảm $>$94\% so với mô hình gốc \\
\hline
Thời gian tải GLB (WiFi 50~Mbps) & 1,2--1,8 giây & Chrome 120, desktop \\
\hline
Thời gian đến tương tác đầu tiên & < 0,8 giây & Quick Tour sẵn sàng trước GLB \\
\hline
Tương thích trình duyệt & Chrome 90+, Firefox 88+, Safari 14+ & Kiểm tra theo bộ kết hợp \\
\hline
\end{tabular}
\end{table}

Giá trị mang lại ở ba chiều. \textit{Với sinh viên:} lần đầu tiên có thể
thực sự ``vào trong'' phòng ký túc xá trước khi đăng ký, nhìn quanh, quan sát
từng chi tiết, cảm nhận tỷ lệ không gian --- thay vì đoán mò qua ảnh tĩnh.
\textit{Với cán bộ quản lý:} công cụ ghi tọa độ hotspot tích hợp trong giao
diện quản trị cho phép cập nhật và mở rộng nội dung trực quan hóa mà không cần
lập trình viên can thiệp. \textit{Về mặt kỹ thuật:} hành trình từ 200~MB về
dưới 8~MB là minh chứng cho thấy các kỹ thuật tối ưu hóa hình học 3D có thể
đạt tỷ lệ nén rất cao khi được áp dụng đúng thứ tự và đúng mức độ cho từng
loại đối tượng trong mô hình.

Nhìn lại, hành trình từ 200~MB về dưới 8~MB dạy tôi điều quan trọng hơn con
số: phải hiểu đúng nguồn gốc vấn đề trước khi bắt đầu tối ưu, không phải cứ
thử từng thứ một.

% =====================================================================
\section{Giải pháp phân bổ phòng tự động dựa trên quota}
\label{sec:contrib_allocation}
% =====================================================================

\subsection{Dẫn dắt và bài toán}
\label{subsec:alloc_problem}

\subsubsection{Vấn đề cần giải quyết}

Bài toán phân bổ phòng ký túc xá mang tính ràng buộc đa chiều: mỗi quyết định
gán phòng phải đồng thời thỏa mãn bốn nhóm ràng buộc --- sức chứa phòng, đồng
nhất giới tính, chỉ tiêu theo nhóm năm học (quota), và thứ tự ưu tiên theo
chính sách nhà trường. Đặc biệt, số lượng đơn đăng ký thường vượt số chỗ ở
khả dụng, đòi hỏi một cơ chế xếp hạng minh bạch và có thể giải trình.

Thực trạng trước khi có hệ thống: quy trình phân bổ được thực hiện thủ công
trên bảng tính điện tử, tiêu tốn nhiều ngày công tác và tiềm ẩn sai sót khi
hợp nhất dữ liệu từ nhiều nguồn. Sinh viên không thể kiểm chứng tại sao mình
được ưu tiên hơn hay kém hơn người khác, dẫn đến mất lòng tin và gia tăng
khiếu nại.

\subsubsection{Đối tượng hưởng lợi}

\textit{Sinh viên đăng ký:} có thể tra cứu lý do cụ thể tại sao mình được
phân phòng hay đưa vào danh sách chờ, thay vì nhận kết quả mà không hiểu cơ
sở. \textit{Cán bộ quản lý:} quy trình phân bổ từ nhiều ngày rút xuống dưới
2 giây, loại bỏ hoàn toàn khả năng sai sót tính điểm thủ công. \textit{Ban
lãnh đạo ký túc xá:} có nhật ký kiểm toán đầy đủ cho mỗi đợt phân bổ, phục
vụ báo cáo và giải trình khi có khiếu nại.

\subsubsection{Điểm khác biệt so với giải pháp hiện có}

\begin{table}[H]
\centering
\caption{So sánh giải pháp phân bổ phòng}
\label{tab:alloc_comparison}
\small
\begin{tabular}{|p{3.5cm}|p{4.5cm}|p{4.8cm}|}
\hline
\textbf{Tiêu chí} & \textbf{Giải pháp thủ công hiện có} & \textbf{Giải pháp đề xuất} \\
\hline
Thời gian xử lý & Nhiều ngày công tác & Dưới 2 giây \\
\hline
Kiểm tra ràng buộc & Dễ bỏ sót khi làm thủ công & Tự động, tường minh, thứ tự cố định \\
\hline
Tính minh bạch & Sinh viên không biết lý do & Mỗi quyết định có AuditLog riêng \\
\hline
Quản lý danh sách chờ & Thủ công, dễ nhầm thứ tự & Tự động, có thứ tự theo điểm \\
\hline
Khả năng kiểm toán & Không có & Ghi đầy đủ căn cứ từng quyết định \\
\hline
\end{tabular}
\end{table}

\subsection{Giải pháp}
\label{subsec:alloc_model}

\subsubsection{Mô hình hai lớp quyết định}

Ban đầu, toàn bộ logic phân bổ được viết trong một hàm duy nhất: kiểm tra
ràng buộc, tính điểm, gán phòng --- tất cả gộp lại. Chỉ sau khi thử test đồng
thời nhiều đơn mới nhận ra vấn đề: khi hai tiến trình cùng đọc số chỗ trống
và cùng gán phòng, kết quả cuối bị sai. Lúc đó mới quyết định tách ``ràng buộc
cứng'' ra kiểm tra trước toàn bộ, và \texttt{capacitySnapshot} xuất hiện như
một bản chụp trạng thái phòng tại thời điểm bắt đầu phân bổ để ngắt race
condition đó.

Mô hình phân bổ hoạt động trên hai lớp quyết định độc lập. \textit{Lớp ngoài}
kiểm tra các ràng buộc cứng theo thứ tự cố định: (1)~quota nhóm năm còn chỗ
không; (2)~phòng còn giường trống không; (3)~giới tính phòng có khớp không.
Chỉ khi thỏa mãn cả ba ràng buộc này, ứng viên mới tiếp tục sang lớp trong.
\textit{Lớp trong} sử dụng điểm ưu tiên để xếp hạng các ứng viên cạnh tranh
trong cùng nhóm:
\begin{equation}
Score_i = W_d \cdot D_i + W_f \cdot F_i + W_y \cdot Y_i + W_p \cdot P_i - W_v \cdot V_i
\label{eq:priority_score}
\end{equation}
trong đó $D_i$, $F_i$, $Y_i$, $P_i$, $V_i$ lần lượt là điểm khoảng cách từ
nhà, tài chính khó khăn, nhóm năm học, học bổng và điểm trừ vi phạm nội quy.
Các trọng số $W_\ast$ được lưu trong cấu hình, có thể điều chỉnh mỗi năm mà
không cần sửa mã nguồn. Điểm ưu tiên là \textit{cơ chế xếp hạng}, không phải
cơ chế quyết định: sinh viên phải vượt qua lớp ngoài trước.

Số chỗ ở dành cho từng nhóm năm được tính theo:
\begin{equation}
Quota_g = Capacity \times Ratio_g
\label{eq:quota}
\end{equation}
với $Ratio_g$ là tỷ lệ được cấu hình trong \texttt{QuotaConfig}. Sau mỗi đợt
phân bổ, tỷ lệ lấp đầy được cập nhật trên bảng điều khiển:
\begin{equation}
OccupancyRate = \frac{OccupiedBeds}{TotalBeds} \times 100\%
\label{eq:occupancy}
\end{equation}
Dữ liệu thực tế ghi nhận 82,1\% trên 265 phòng và 1.450 sinh viên.

\subsubsection{Thuật toán phân bổ phòng}

\begin{minipage}{\linewidth}
\paragraph{Mục tiêu.}
Phân bổ toàn bộ danh sách đơn đăng ký vào các phòng khả dụng theo đúng thứ
tự ưu tiên, đảm bảo ba ràng buộc cứng không bao giờ bị vi phạm và mọi quyết
định đều có căn cứ ghi lại trong AuditLog.

\begin{algorithm}[H]
\small
\caption{\textit{QuotaBasedRoomAllocation}}
\KwData{$StudentList$: danh sách đơn; $RoomList$: danh sách phòng; $QuotaConfig$: cấu hình quota}
\KwResult{$AssignmentList$: danh sách phân phòng; $WaitingList$: danh sách chờ}
\tcp{Bước 1: Tính điểm và sắp xếp}
\For{$s \in StudentList$}{
  $s.score \leftarrow \text{CalculatePriorityScore}(s)$\;
}
$StudentList \leftarrow \text{SORT}(StudentList,\ key{:}\ score,\ \text{DESC})$\;
\tcp{Bước 2: Phân bổ tuần tự --- kiểm tra 3 ràng buộc cứng}
\For{$s \in StudentList$}{
  \If{$QuotaConfig.slots[s.yearGroup] \leq 0$}{
    $WaitingList.\text{append}(s,\ reason{:}\ \text{QUOTA\_FULL})$\;
    \KwBreak\;
  }
  $room \leftarrow \text{FindRoom}(RoomList,\ gender{:}\ s.gender,\ group{:}\ s.yearGroup)$\;
  \If{$room \neq \varnothing$}{
    $\text{Assign}(s,\ room)$\;
    $QuotaConfig.slots[s.yearGroup] \mathrel{-}= 1$\;
    $\text{WriteAuditLog}(s,\ room,\ \text{ALLOCATED})$\;
  }
  \Else{
    $WaitingList.\text{append}(s,\ reason{:}\ \text{NO\_ROOM})$\;
    $\text{WriteAuditLog}(s,\ \text{WAITLIST})$\;
  }
}
\KwReturn $AssignmentList$, $WaitingList$\;
\end{algorithm}
\end{minipage}

\medskip
\paragraph{Phân tích độ phức tạp.}
Bước tính điểm là $O(n)$; sắp xếp là $O(n \log n)$; phân bổ là $O(n \cdot m)$
trong trường hợp xấu nhất, giảm về $O(n \cdot m/k)$ khi danh sách phòng được
lập chỉ mục theo nhóm năm và giới tính. Với 1.500 đơn và 265 phòng, thuật toán
hoàn thành dưới 2 giây. Tính đúng đắn được đảm bảo bởi $capacitySnapshot$ bất
biến: bản chụp sức chứa ghi lại trước khi phân bổ bắt đầu, tránh race condition
khi cán bộ thao tác phòng đồng thời.

\subsection{Kết quả đạt được}
\label{subsec:alloc_result}

\begin{table}[H]
\centering
\caption{Kết quả kiểm thử xác minh tính đúng đắn thuật toán phân bổ}
\label{tab:alloc_verification}
\small
\begin{tabular}{|p{3.0cm}|p{3.6cm}|p{3.8cm}|p{1.8cm}|}
\hline
\textbf{Kịch bản} & \textbf{Điều kiện đầu vào} & \textbf{Kết quả mong đợi} & \textbf{Kết quả} \\
\hline
Đủ sức chứa & 10 đơn, 15 giường trống & Tất cả 10 đơn ALLOCATED & Đạt \\
\hline
Vượt sức chứa & 20 đơn, 10 giường & 10 điểm cao nhất ALLOCATED, 10 WAITLIST & Đạt \\
\hline
Xếp hạng đúng & 5 đơn điểm khác nhau, 3 giường & 3 đơn điểm cao nhất theo thứ tự giảm dần & Đạt \\
\hline
Hết quota nhóm & 10 đơn năm 1, quota = 5 & 5 ALLOCATED, 5 QUOTA\_FULL & Đạt \\
\hline
Không có giường & Bất kỳ số đơn, 0 giường & Tất cả NO\_ROOM & Đạt \\
\hline
\end{tabular}
\end{table}

So với quy trình thủ công, hệ thống rút ngắn thời gian phân bổ từ vài ngày
xuống dưới 2 giây, loại bỏ sai số tính điểm tay, và đảm bảo mọi sinh viên
đều có thể tra cứu căn cứ quyết định của mình trong \texttt{AllocationAuditLog}.
Ba ràng buộc cứng được kiểm tra tuần tự và tường minh trong mã nguồn, không
thể bị bỏ qua dù có lỗi ở bước khác.

Điều đọng lại không phải là tốc độ dưới 2 giây, mà là lần đầu tiên sinh viên
có thể tự tra được lý do tại sao mình được hay không được phân phòng ---
minh bạch đó mới là giá trị thực.

% =====================================================================
\section{Giải pháp kiến trúc phân bổ linh hoạt theo năm học}
\label{sec:contrib_policy}
% =====================================================================

\subsection{Dẫn dắt và bài toán}
\label{subsec:policy_problem}

\subsubsection{Vấn đề cần giải quyết}

Chính sách phân bổ ký túc xá không cố định mà thay đổi mỗi năm học: tỷ lệ
ưu tiên theo nhóm năm có thể điều chỉnh, tiêu chí xét duyệt có thể mở rộng,
ngưỡng khoảng cách có thể thay đổi. Trong hệ thống với logic viết cứng trong
mã nguồn, mỗi thay đổi chính sách đòi hỏi một chu kỳ phát triển phần mềm mới
--- độ trễ hoàn toàn không phù hợp với nhịp thay đổi năm học.

Bên cạnh đó, cơ cấu sinh viên theo nhóm năm dịch chuyển mỗi năm (cohort
shift): khi K69 lên năm 2, K70 vào thay vị trí năm 1. Nếu không có công cụ
dự báo, cán bộ buộc phải phân bổ chỉ tiêu dựa trên kinh nghiệm chủ quan thay
vì dữ liệu thực tế. Đây là bài toán lập kế hoạch chủ động, không phải phản
ứng sau khi vấn đề đã xảy ra.

\subsubsection{Đối tượng hưởng lợi}

\textit{Cán bộ quản lý:} có thể tự cập nhật chính sách phân bổ hàng năm qua
giao diện cấu hình mà không cần yêu cầu phát triển phần mềm mới. \textit{Ban
lãnh đạo ký túc xá:} có lịch sử chính sách đầy đủ theo từng năm học để đối
chiếu và kiểm toán. \textit{Toàn bộ sinh viên:} được hưởng lợi từ chính sách
phân bổ luôn phản ánh đúng thực tế cơ cấu năm học hiện tại.

\subsubsection{Điểm khác biệt so với giải pháp hiện có}

\begin{table}[H]
\centering
\caption{So sánh giải pháp quản lý chính sách phân bổ}
\label{tab:policy_comparison}
\small
\begin{tabular}{|p{3.5cm}|p{4.5cm}|p{4.8cm}|}
\hline
\textbf{Tiêu chí} & \textbf{Logic hard-coded hiện có} & \textbf{Policy-Driven Framework} \\
\hline
Thay đổi chính sách & Cần chu kỳ phát triển phần mềm & Cấu hình qua giao diện, vài phút \\
\hline
Kế thừa năm trước & Không có, cấu hình lại từ đầu & Tự động kế thừa, admin chỉ điều chỉnh \\
\hline
Lịch sử chính sách & Không được lưu trữ & Lưu đầy đủ theo phiên bản và năm học \\
\hline
Dự báo cohort shift & Không có & Module tự động dựa trên dữ liệu lịch sử \\
\hline
Tách biệt quota và điểm & Không (lẫn lộn trong code) & Tách thành hai thực thể độc lập \\
\hline
\end{tabular}
\end{table}

\subsection{Giải pháp}
\label{subsec:policy_arch}

\subsubsection{Kiến trúc bốn lớp}

Lúc đầu, \texttt{QuotaConfig} được lồng trực tiếp bên trong
\texttt{AllocationPolicy} --- nghe có vẻ hợp lý vì chúng liên quan đến nhau.
Nhưng sau một lần cán bộ yêu cầu điều chỉnh tỷ lệ quota mà không muốn đụng
vào công thức tính điểm ưu tiên, mới nhận ra hai thứ này có vòng đời hoàn
toàn khác nhau: quota thay đổi theo cơ cấu sinh viên từng năm, còn trọng số
điểm ưu tiên có thể giữ nguyên nhiều năm liền. Đó là lúc kiến trúc bốn lớp
hình thành.

Hình~\ref{fig:policy_driven_allocation} trình bày kiến trúc Policy-Driven
Allocation Framework: chính sách phân bổ được lưu trữ dưới dạng tài liệu có
cấu trúc trong cơ sở dữ liệu, hoàn toàn tách khỏi mã nguồn.

\begin{figure}[H]
\centering
\includegraphics[width=0.85\textwidth]{Hinhve/policy_driven_allocation.png}
\caption{Kiến trúc quản lý chính sách phân bổ động}
\label{fig:policy_driven_allocation}
\end{figure}

Kiến trúc gồm bốn lớp tương tác theo chiều dọc. \texttt{QuotaConfig} định
nghĩa tỷ lệ chỉ tiêu theo nhóm năm --- thực thể thay đổi hàng năm theo biến
động cơ cấu sinh viên. \texttt{AllocationPolicy} định nghĩa trọng số tính
điểm ưu tiên và tham chiếu về \texttt{QuotaConfig} --- thực thể ổn định hơn,
có thể dùng qua nhiều năm. \texttt{AllocationCycle} xác định khoảng thời gian
đăng ký và liên kết với một chính sách cụ thể. \texttt{AllocationService} đọc
cấu hình từ ba lớp trên và thực thi phân bổ mà không chứa bất kỳ giá trị
hard-coded nào.

Quyết định thiết kế then chốt là tách \texttt{QuotaConfig} khỏi
\texttt{AllocationPolicy}: hai khái niệm nghiệp vụ có vòng đời thay đổi khác
nhau không nên lưu trong cùng một cấu trúc dữ liệu. Nhờ tách biệt này, cán bộ
có thể điều chỉnh tỷ lệ chỉ tiêu mà không ảnh hưởng đến công thức tính điểm
ưu tiên, và ngược lại.

Công thức theo dõi mức độ tuân thủ quota sau mỗi đợt phân bổ:
\begin{equation}
\Delta_g = n_g^{\mathrm{thực\ tế}} - \lfloor Q_g \cdot C_{\mathrm{tổng}} \rfloor
\label{eq:quota_deviation}
\end{equation}
Khi $\Delta_g > 0$, nhóm $g$ đang vượt chỉ tiêu; $\Delta_g < 0$ là chưa đạt.
Hệ thống hiển thị chỉ số này theo thời gian thực để cán bộ điều chỉnh cho các
đợt đăng ký tiếp theo.

\subsubsection{Thuật toán kế thừa chính sách}

\begin{minipage}{\linewidth}
\paragraph{Mục tiêu.}
Đảm bảo chu kỳ đăng ký mới luôn có chính sách hợp lệ mà cán bộ không phải
cấu hình lại từ đầu mỗi năm, đồng thời bảo toàn toàn bộ lịch sử chính sách
theo phiên bản.

\begin{algorithm}[H]
\small
\caption{\textit{PolicyInheritanceAndActivation}}
\KwData{$academicYear$: năm học mới; $adminRequest$: yêu cầu tạo chu kỳ}
\KwResult{$newPolicy$: chính sách hợp lệ; $cycle$: chu kỳ đăng ký đã liên kết}
\tcp{Ưu tiên 1: Tìm chính sách đang hiệu lực}
$activePolicy \leftarrow \text{FindActivePolicy}(academicYear)$\;
\If{$activePolicy \neq \varnothing$}{
  $newPolicy \leftarrow activePolicy$\;
}
\Else{
  \tcp{Ưu tiên 2: Kế thừa từ năm trước, chờ admin xét duyệt}
  $prevPolicy \leftarrow \text{FindActivePolicy}(\text{PreviousYear}(academicYear))$\;
  \If{$prevPolicy \neq \varnothing$}{
    $newPolicy \leftarrow \text{ClonePolicy}(prevPolicy,\ academicYear)$\;
    $newPolicy.active \leftarrow \text{FALSE}$\;
  }
  \Else{
    \tcp{Ưu tiên 3: Tạo chính sách mặc định nếu không có tiền lệ}
    $newPolicy \leftarrow \text{CreateDefaultPolicy}(academicYear)$\;
  }
  $\text{SavePolicy}(newPolicy)$\;
}
\tcp{Liên kết chu kỳ với chính sách}
$cycle \leftarrow \text{CreateCycle}(adminRequest,\ policyId{:}\ newPolicy.\_id)$\;
$cycle.capacitySnapshot \leftarrow \text{CaptureRoomCapacity}()$\;
\KwReturn $newPolicy$, $cycle$\;
\end{algorithm}
\end{minipage}

\medskip
\paragraph{Phân tích độ phức tạp.}
Thuật toán có độ phức tạp $O(r)$ với $r$ là số phòng (bước CaptureRoomCapacity);
ClonePolicy là $O(1)$ với cấu trúc policy kích thước cố định. Với 265 phòng,
toàn bộ thuật toán hoàn thành trong vài trăm mili giây. Phân cấp ưu tiên ba
bậc đảm bảo luôn có đầu ra hợp lệ trong mọi tình huống, kể cả năm học đầu
tiên chưa có tiền lệ.

\subsection{Kết quả đạt được}
\label{subsec:policy_result}

Kiến trúc Policy-Driven Allocation mang lại ba giá trị định lượng được.
\textit{Thứ nhất}, thời gian cập nhật chính sách mỗi năm giảm từ vài ngày
(phát triển phần mềm mới) xuống còn vài phút (điền form cấu hình). \textit{Thứ
hai}, lịch sử chính sách được bảo toàn đầy đủ theo năm học thông qua trường
\texttt{version} và \texttt{sourceQuotaId}, hỗ trợ kiểm toán và đối chiếu.
\textit{Thứ ba}, module dự báo cohort shift đề xuất quota tự động dựa trên
xu hướng $k$ năm gần nhất:
\begin{equation}
\hat{Q}_g^{\mathrm{đề\ xuất}} = \frac{1}{k} \sum_{i=1}^{k}
    \frac{n_{g,i}^{\mathrm{thực\ tế}}}{\displaystyle\sum_{g'} n_{g',i}^{\mathrm{thực\ tế}}}
\label{eq:quota_recommendation}
\end{equation}
cho phép cán bộ lập kế hoạch chủ động trước khi năm học mới bắt đầu thay vì
phản ứng sau khi vấn đề đã xảy ra.

Tách cấu hình khỏi code không chỉ là quyết định kỹ thuật --- nó là quyết định
về ai có quyền thay đổi hệ thống mà không cần lập trình viên.

% =====================================================================
\section{Tổng kết chương}
\label{sec:contrib_summary}
% =====================================================================

Chương 5 trình bày ba đóng góp kỹ thuật của hệ thống eDorm từ góc nhìn tư duy
thiết kế và giải quyết vấn đề. Mỗi đóng góp được đặt trong bối cảnh vấn đề
thực tế quan sát được, phân tích điểm khác biệt so với giải pháp hiện có, và
đánh giá giá trị mang lại cho từng nhóm đối tượng.

\textbf{Đóng góp 1} --- trực quan hóa không gian phòng
(mục~\ref{sec:contrib_visualization}) --- giải quyết bài toán sinh viên không
thể đánh giá phòng ở trước khi đăng ký. Hành trình bốn giai đoạn từ ảnh tĩnh
đến kiến trúc hai chế độ song song (Quick Tour 360° và Free Explore 3D), đặc
biệt là quá trình tối ưu hóa mô hình 3D từ 200~MB về dưới 8~MB, phản ánh quá
trình tinh chỉnh liên tục dựa trên những giới hạn thực tế gặp phải. Điểm đặc
trưng là kết hợp có chủ ý hai công nghệ theo điểm mạnh của từng loại thiết bị,
thay vì áp dụng một giải pháp đồng nhất.

\textbf{Đóng góp 2} --- phân bổ phòng tự động dựa trên quota
(mục~\ref{sec:contrib_allocation}) --- chuyển hóa quy trình nhiều ngày thủ
công thành tự động dưới 2 giây. Điểm cốt lõi là mô hình hai lớp quyết định
tách biệt ràng buộc cứng (quota, giới tính, sức chứa) khỏi cơ chế xếp hạng
(điểm ưu tiên), kết hợp với AuditLog đảm bảo tính minh bạch và khả năng giải
trình cho từng sinh viên.

\textbf{Đóng góp 3} --- kiến trúc phân bổ linh hoạt theo năm học
(mục~\ref{sec:contrib_policy}) --- giải quyết bài toán dài hạn: làm sao chính
sách phân bổ có thể thay đổi mỗi năm mà không cần sửa mã nguồn. Kiến trúc
bốn lớp tách hoàn toàn cấu hình khỏi logic thực thi, kết hợp cơ chế kế thừa
chính sách và module dự báo cohort shift để hỗ trợ lập kế hoạch chủ động.

\end{document}
